% frontmatter/actividades.tex -- las actividades, una a una, para el
% adulto: cada estación trae las suyas (ver notes/01-plan.md), y esta
% página tiene que seguir cabiendo en una (tools/check_pages.py comprueba
% que el día 1 está en la página 5).
\section*{Las actividades}

Cada actividad llega el día que se usa por primera vez, y después vuelve
de vez en cuando. Las instrucciones de cada página, en letra pequeña y
gris, las lee un adulto.

{\small
\begin{multicols}{2}
\raggedright
\subsection*{Desde el otoño}
\begin{itemize}[leftmargin=5mm, itemsep=1pt, topsep=0pt]
\item \textbf{\lblTraza} (los lunes, cuando llega un número nuevo):
  repasar el número con el dedo por encima de los puntos, luego con un
  lápiz, y después escribirlo solo en la línea. Los puntos son el
  contorno del número de verdad, el mismo que se lee en el resto del
  cuaderno.
\item \textbf{\lblColorea}: colorear exactamente las que dice, ni una
  más -- las demás se quedan en blanco.
\item \textbf{\lblRodea}: entre varios grupos de cosas, rodear el que
  tiene las que dice. Hasta el 6, los grupos están puestos como los
  puntos de un dado, y del 7 al 10 como en el marco de diez, para que se
  vean de un vistazo; el grupo vacío es el 0.
\item \textbf{\lblCuenta}: contar señalando cada cosa con el dedo, una
  sola vez cada una, y rodear el número.
\item \textbf{\lblBusca}: rodear el número cada vez que aparece, entre
  otros números y formas.
\item \textbf{\lblUne}: contar cada grupo y trazar una línea desde su
  punto hasta el punto de su número.
\item \textbf{\lblCompleta}: decir en voz alta los números del tren, uno
  a uno, y escribir en los vagones vacíos los que faltan.
\item \textbf{\lblDibuja}: dibujar las cosas que dice, ni más ni menos.
\item \textbf{\lblRepasa} (los viernes): lo que ya sabe hacer esa
  semana, para marcarlo, y un dibujo. Cada diez páginas hay un pequeño
  cartel de celebración.
\end{itemize}

\subsection*{Desde el invierno}
\begin{itemize}[leftmargin=5mm, itemsep=1pt, topsep=0pt]
\item \textbf{\lblDecena} (del 11 al 20): una bandeja llena, que son 10,
  y otra al lado; se sigue contando desde el 10.
\item \textbf{\lblCompara}: rodear el grupo que tiene más, o menos, o
  los dos que tienen las mismas; o el número más grande, o el más
  pequeño.
\item \textbf{\lblVecinos}: escribir el número que va antes y el que va
  después, en las casas de los lados.
\item \textbf{\lblRecta}: decir los números de la recta, de izquierda a
  derecha, y escribir los que faltan.
\item \textbf{\lblOrdena}: escribir los números en orden, del más
  pequeño al más grande (o al revés).
\item \textbf{\lblJunta}: juntar dos grupos y contar cuántas son en
  total: la suma, antes de los signos + e =.
\end{itemize}

\subsection*{Desde la primavera}
\begin{itemize}[leftmargin=5mm, itemsep=1pt, topsep=0pt]
\item \textbf{\lblSuma} y \textbf{\lblResta}: lo mismo, ya con sus
  signos; en la resta, se tachan las que se quitan.
\item \textbf{\lblParte}: una raya parte el grupo en dos; se cuenta lo
  que queda al otro lado.
\item \textbf{\lblDiez}: dibujar los puntos que faltan para llenar el
  marco de diez.
\item \textbf{\lblProblema}: lo lee el adulto; se puede dibujar, y se
  escribe la cuenta.
\item \textbf{\lblBloques}: cada barra es una decena (diez cubitos);
  se cuentan las barras y los cubitos sueltos.
\item \textbf{\lblCompleta}, \textbf{\lblRecta} y \textbf{\lblCuenta},
  también de 2 en 2, de 5 en 5 y de 10 en 10: las ruedas de las bicis,
  los dedos de las manos, las decenas.
\end{itemize}
\end{multicols}}
\newpage
